\begin{figure}[H]
    \centering
    \foreach \l in {0,...,10}
    {
      \includegraphics[width=4.2cm]{linewidth/pattern\l-45.png}%
    }%
    \caption{
      Line Mode with Pulse Width at Maximum
      }
\end{figure}
\clearpage
\section*{line mode} 
\label{sec:linemode}
\lstset{style=6502Style}

\begin{definition}[Jeffrey Says]
\setlength{\intextsep}{0pt}%
\setlength{\columnsep}{3pt}%
\begin{wrapfigure}{l}{0.12\textwidth}
\includegraphics[width=\linewidth]{src/callout/psych.png} 
\end{wrapfigure}
\textbf{Press L to turn on and off} the Line Mode - a bit like drawing with the Aurora Borealis.\\
\textbf{Press W to adjust line width:} Sets the width of the lines produced in Line Mode.
\\
\\
\end{definition}
'Line Mode' is a completely different mode of painting. 

\clearpage
\textbf{Lines 1189-1231. \icode{\textbf{JustLPressed}}} 
\begin{lstlisting}[caption=From \icode{CheckKeyboardInput}.]
JustLPressed   
        ; 'L' pressed. Turn line mode on or off.
        LDA lineModeActivated
        EOR #$01
        STA lineModeActivated
        ASL 
        ASL 
        ASL 
        ASL 
        TAY 

        ; Briefly display the new linemode on the bottom of the screen.
        JSR ClearLastLineOfScreen
        LDX #$00
_Loop   LDA lineModeSettingDescriptions,Y
        STA lastLineBufferPtr,X
        INY 
        INX 
        CPX #$10
        BNE _Loop
        JMP WriteLastLineBufferToScreen
        ; Returns
\end{lstlisting}

\clearpage

\textbf{Lines 1189-1231. \icode{\textbf{JustLPressed}}:} 
Pressing \icode{L} toggles line mode on or off. Line mode's on/off state is 
stored in \icode{lineModeActivated}, with \icode{\$00} meaning 'off' and
\icode{\$01} meaning 'on'. This means we can use a neat, economical trick
for toggling the value every time the user presses the 'L' key:

\begin{lstlisting}
        LDA lineModeActivated
        EOR #$01
        STA lineModeActivated
\end{lstlisting}

The \icode{EOR} statement performs an exclusive-or that has the neat property of
turning \icode{\$00} into \icode{\$01} and \icode{\$01} into \icode{\$00} - equivelent
to switching a value on or off.

This is what the the bit by bit operation looks like when '\icode{EOR \$01}' turns
\icode{trackingActivated} 'on':

\begin{figure}[H]
  {
    \setlength{\tabcolsep}{3.0pt}
    \setlength\cmidrulewidth{\heavyrulewidth} % Make cmidrule = 
    \begin{adjustbox}{width=7cm,center}

      \begin{tabular}{rllllllll}
        \toprule
        Byte & Bit 7 & Bit 6 & Bit 5 & Bit 4 & Bit 3 & Bit 2 & Bit 1 & Bit 0        \\
        \midrule
        \$00 & 0 & 0 & 0 & 0 & 0 & 0 & 0 & 0 \\
        \$01 & 0 & 0 & 0 & 0 & 0 & 0 & 0 & 1 \\
        \midrule
        Result & 0 & 0 & 0 & 0 & 0 & 0 & 0 & 1 \\
        \addlinespace
        \bottomrule
      \end{tabular}

    \end{adjustbox}

  }\caption*{X-OR'ing \$01 and \$00 gives \$01, the 'on' value for \icode{lineModeActivated}.}
\end{figure}

And this is what it looks like when \icode{EOR \$01} turns \icode{lineModeActivated} 'off' again:
\begin{figure}[H]
  {
    \setlength{\tabcolsep}{3.0pt}
    \setlength\cmidrulewidth{\heavyrulewidth} % Make cmidrule = 
    \begin{adjustbox}{width=7cm,center}

      \begin{tabular}{rllllllll}
        \toprule
        Byte & Bit 7 & Bit 6 & Bit 5 & Bit 4 & Bit 3 & Bit 2 & Bit 1 & Bit 0        \\
        \midrule
        \$01 & 0 & 0 & 0 & 0 & 0 & 0 & 0 & 1 \\
        \$01 & 0 & 0 & 0 & 0 & 0 & 0 & 0 & 1 \\
        \midrule
        Result & 0 & 0 & 0 & 0 & 0 & 0 & 0 & 0 \\
        \addlinespace
        \bottomrule
      \end{tabular}

    \end{adjustbox}

  }\caption*{X-OR'ing \$01 and \$01 gives \$00, the 'off' value for \icode{lineModeActivated}.}
\end{figure}
\clearpage

\textbf{Lines 1189-1231. \icode{\textbf{ApplyLineMode}}} 
\begin{lstlisting}[caption=From \icode{MainInterruptHandler}.]
;-------------------------------------------------------
; MainInterruptHandler
;-------------------------------------------------------
MainInterruptHandler
        ...
ApplyLineMode
        LDA lineModeActivated
        BEQ LineModeNotActive

        ; Line Mode Active
        LDA #NUM_ROWS + 1
        SEC 
        SBC cursorYPosition
        ORA #$80
        STA currentColorIndexArray,X
\end{lstlisting}

\begin{lstlisting}[caption=From \icode{MainPaintLoop}.]
;-------------------------------------------------------
; MainPaintLoop
;-------------------------------------------------------
MainPaintLoop    
        ...
        ; Line Mode sets the top bit of currentValueInColorIndexArray
        LDA currentValueInColorIndexArray
        AND #$80   ; #LINE_MODE_ACTIVE
        BNE PaintLineModeAndLoop
        ...
PaintLineModeAndLoop
        ; Loops back to MainPaintLoop
        JMP PaintLineMode
\end{lstlisting}
\clearpage
\textbf{Lines 1189-1231. \icode{\textbf{ApplyLineMode}}:} Here \icode{lineModeActivated} is used
to tag one of the pixel buffers to indicate that line mode, rather than the usual painting routine
in \icode{PaintStructureAtCurrentPosition} should be used for painting. This 'tagging'
is done by setting the leftmost bit of the byte stored to the current position in the 
\icode{currentColorIndexArray}. We do this by 'OR'ing the manipulated Y position of the cursor
with \icode{\$80} before we store it:

\begin{figure}[H]
  {
    \setlength{\tabcolsep}{3.0pt}
    \setlength\cmidrulewidth{\heavyrulewidth} % Make cmidrule = 
    \begin{adjustbox}{width=7cm,center}

      \begin{tabular}{rllllllll}
        \toprule
        Byte & Bit 7 & Bit 6 & Bit 5 & Bit 4 & Bit 3 & Bit 2 & Bit 1 & Bit 0        \\
        \midrule
        \$15 & 0 & 0 & 0 & 1 & 0 & 1 & 0 & 1 \\
        \$80 & 1 & 0 & 0 & 0 & 0 & 0 & 0 & 0 \\
        \midrule
      Result & 1 & 0 & 0 & 1 & 0 & 1 & 0 & 1 \\
        \addlinespace
        \bottomrule
      \end{tabular}

    \end{adjustbox}

  }\caption*{OR'ing a \icode{cursorYPosition} of \$15 and \$80 gives \$95, a value of \icode{\$15} with the leftmost bit set.}
\end{figure}

This is not the only strange thing happening here, it also 
appears we storing a Y co-ordinate in the color index array, an array we typically use for
storing color values. Specifically we're taking the 'reflected' vertical position of the 
cursor and storing it there. Clearly line mode is a different beast from our normal painting routine. We will
see what lies behind this a little later.


\textbf{Lines 1189-1231. \icode{\textbf{MainPaintLoop}}:} For now, in the main paint loop, when it comes time to do some
painting, we use the 'tag' we set above to detect if we should paint in 'Line Mode'. We do this by \icode{AND}'ing the
value we loaded from \icode{currentColorIndexArray} with \icode{\$80} to detect if it has been tagged:

\begin{figure}[H]
  {
    \setlength{\tabcolsep}{3.0pt}
    \setlength\cmidrulewidth{\heavyrulewidth} % Make cmidrule = 
    \begin{adjustbox}{width=7cm,center}

      \begin{tabular}{rllllllll}
        \toprule
        Byte & Bit 7 & Bit 6 & Bit 5 & Bit 4 & Bit 3 & Bit 2 & Bit 1 & Bit 0        \\
        \midrule
        \$95 & 1 & 0 & 0 & 1 & 0 & 1 & 0 & 1 \\
        \$80 & 1 & 0 & 0 & 0 & 0 & 0 & 0 & 0 \\
        \midrule
      Result & 1 & 0 & 0 & 0 & 0 & 0 & 0 & 0 \\
        \addlinespace
        \bottomrule
      \end{tabular}

    \end{adjustbox}

  }\caption*{AND'ing a \icode{currentValueInColorIndexArray} of \$95 and \$80 gives \$80 - this non-zero value tells us Line Mode is active.}
\end{figure}

If line mode is active, we \icode{JMP} to \icode{PaintLineMode} which we cover next.

\clearpage
\textbf{Lines 1189-1231. \icode{\textbf{PaintLineMode}}} 
\begin{lstlisting}[caption=From \icode{PaintLineMode}.]
PaintLineMode 
        LDA currentValueInColorIndexArray
        AND #$7F
        STA offsetForYPos
        LDA #NUM_ROWS + 1
        SEC 
        SBC offsetForYPos
        STA pixelYPosition
        DEC pixelYPosition
        LDA #$00
        STA currentValueInColorIndexArray
        LDA #ACTIVE
        STA skipPixel
        JSR PaintPixelForCurrentSymmetry
        INC pixelYPosition
        LDA #NOT_ACTIVE
        STA skipPixel

        LDA lineWidth
        EOR #$07
        STA currentValueInColorIndexArray
LineModeLoop   
        JSR PaintPixelForCurrentSymmetry
        INC pixelYPosition
        INC currentValueInColorIndexArray
        LDA currentValueInColorIndexArray
        CMP #$08
        BNE ResetLineModeColorValue
        JMP CleanUpAndExitLineModePaint

        INC currentValueInColorIndexArray
ResetLineModeColorValue   
        STA currentValueInColorIndexArray
        LDA pixelYPosition
        CMP #NUM_ROWS + 1
        BNE LineModeLoop

CleanUpAndExitLineModePaint    
        LDX currentIndexToPixelBuffers
        DEC currentColorIndexArray,X
        LDA currentColorIndexArray,X
        CMP #$80
        BEQ ResetIndexAndExitLineModePaint
        JMP MainPaintLoop

ResetIndexAndExitLineModePaint   
        LDA #$FF
        STA currentColorIndexArray,X
        STX previousIndexToPixelBuffers
        JMP MainPaintLoop
\end{lstlisting}
\clearpage

\textbf{Lines 1189-1231. \icode{\textbf{PaintLineMode}}:} 
\clearpage
